\documentclass[11pt]{article}
% Shared preamble for per-Python-file documentation (input via \input{...}).
\usepackage[margin=1in]{geometry}
\usepackage[T1]{fontenc}
\usepackage[utf8]{inputenc}
\usepackage{lmodern}
\usepackage{hyperref}
\usepackage{xcolor}
\usepackage{listings}
\usepackage{listingsutf8}
\lstset{
  basicstyle=\ttfamily\small,
  columns=fullflexible,
  breaklines=true,
  upquote=true,
  inputencoding=utf8,
  extendedchars=true,
  frame=single,
  framerule=0.2pt,
  tabsize=2
}


\title{\texttt{dataloader/\_\_init\_\_.py} (Q\&A)}
\author{}
\date{}

\begin{document}
\maketitle

\paragraph{Q: What is the purpose of \texttt{dataloader/\_\_init\_\_.py}?}
\textbf{A:} It makes the dataloader package easier to import by re-exporting the main loader class \lstinline|KUAIRECDataLoader| at the package level, so other scripts can do \lstinline|from dataloader import KUAIRECDataLoader|.

\paragraph{Q: What symbols does it expose?}
\textbf{A:} It imports \lstinline|KUAIRECDataLoader| from \lstinline|dataloader/kuairec.py| and exposes it as \lstinline|dataloader.KUAIRECDataLoader|.

\paragraph{Q: What is an example of how other files use it?}
\textbf{A:} Run scripts (e.g. \lstinline|run_egmn.py|) rely on this re-export to keep their imports short:
\begin{lstlisting}[language=Python]
from dataloader import KUAIRECDataLoader
\end{lstlisting}

\paragraph{Q: Full source code?}
\textbf{A:}
\lstinputlisting[language=Python]{/workspace/dataloader/__init__.py}

\end{document}
